\section{Finite-prefix base and explicit orbit expansion}

This chapter records the finite-prefix base in full.  The base is the
only finite data used by the proof.  Every entry is an exact
arithmetic identity, expanded by rewriting in Lean rather than by
\texttt{native\_decide}.

\subsection{The 17 states}

The accelerated $5x+1$ orbit of $7$ begins with
\begin{equation*}
  7,9,23,29,73,183,229,573,1433,3583,4479,5599,6999,8749,
  21873,54683,34177.
\end{equation*}
There are seventeen entries, indexed by $n=0,\dots,16$.  The proof
never needs a later state.

\subsection{Exact step table}

Each row below records the exact equation
\begin{equation*}
  5\,\Orbit_7(n)+1=2^{t_n}\,\Orbit_7(n+1),
\end{equation*}
with $t_n=\vtwo{5\,\Orbit_7(n)+1}$.

\begin{center}
\scriptsize
\begin{tabular}{@{}rrrrr@{}}
\toprule
$n$ & $\Orbit_7(n)$ & $5\Orbit_7(n)+1$ & $t_n$ & $\Orbit_7(n+1)$ \\
\midrule
0 & 7 & 36 & 2 & 9 \\
1 & 9 & 46 & 1 & 23 \\
2 & 23 & 116 & 2 & 29 \\
3 & 29 & 146 & 1 & 73 \\
4 & 73 & 366 & 1 & 183 \\
5 & 183 & 916 & 2 & 229 \\
6 & 229 & 1146 & 1 & 573 \\
7 & 573 & 2866 & 1 & 1433 \\
8 & 1433 & 7166 & 1 & 3583 \\
9 & 3583 & 17916 & 2 & 4479 \\
10 & 4479 & 22396 & 2 & 5599 \\
11 & 5599 & 27996 & 2 & 6999 \\
12 & 6999 & 34996 & 2 & 8749 \\
13 & 8749 & 43746 & 1 & 21873 \\
14 & 21873 & 109366 & 1 & 54683 \\
15 & 54683 & 273416 & 3 & 34177 \\
\bottomrule
\end{tabular}
\end{center}

The divisibility checks are literal:
\begin{align*}
  36&=4\cdot9, & 46&=2\cdot23, & 116&=4\cdot29, \\
  146&=2\cdot73, & 366&=2\cdot183, & 916&=4\cdot229, \\
  1146&=2\cdot573, & 2866&=2\cdot1433, & 7166&=2\cdot3583, \\
  17916&=4\cdot4479, & 22396&=4\cdot5599, & 27996&=4\cdot6999, \\
  34996&=4\cdot8749, & 43746&=2\cdot21873, & 109366&=2\cdot54683, \\
  273416&=8\cdot34177.
\end{align*}

\subsection{Step weights}

Reading the third column from the table gives the exact weight word
\begin{equation*}
  t_0,\dots,t_{15}
  =2,1,2,1,1,2,1,1,1,2,2,2,2,1,1,3.
\end{equation*}
The first fifteen weights, for $n<15$, are all at most $2$.  The first
weight of size at least $3$ occurs at $n=15$ and equals $3$ exactly.

\begin{lemma}[First big step]\label{sup-lem:first-big}
For every $n<15$, $t_n\le2$, and
\begin{equation*}
  t_{15}=3.
\end{equation*}
Consequently the first step of weight at least $3$ is unique and is
the step from depth $15$ to depth $16$.
\end{lemma}

\begin{proof}
The first statement is the inspection of the table above.  The
uniqueness follows because all earlier weights are at most $2$ and the
weight at depth $15$ is $3$.
\end{proof}

\subsection{Prefix ledger}

Let $\wordWeight_n$ be the prefix weight after $n$ steps and let
$\wordA_n$ be the word molecule of the exact prefix.  The table below
repeats the main-paper ledger with the exact prefix word displayed.

\begin{center}
\scriptsize
\begin{tabular}{@{}rlrrr@{}}
\toprule
$n$ & prefix word & $t_n$ & $\wordWeight_n$ & $\wordA_n$ \\
\midrule
0 & $[]$ & -- & 0 & 0 \\
1 & $[2]$ & 2 & 2 & 1 \\
2 & $[2,1]$ & 1 & 3 & 9 \\
3 & $[2,1,2]$ & 2 & 5 & 53 \\
4 & $[2,1,2,1]$ & 1 & 6 & 297 \\
5 & $[2,1,2,1,1]$ & 1 & 7 & 1549 \\
6 & $[2,1,2,1,1,2]$ & 2 & 9 & 7873 \\
7 & $[2,1,2,1,1,2,1]$ & 1 & 10 & 39877 \\
8 & $[2,1,2,1,1,2,1,1]$ & 1 & 11 & 200409 \\
9 & $[2,1,2,1,1,2,1,1,1]$ & 1 & 12 & 1004093 \\
10 & $[2,1,2,1,1,2,1,1,1,2]$ & 2 & 14 & 5024561 \\
11 & $[2,1,2,1,1,2,1,1,1,2,2]$ & 2 & 16 & 25139189 \\
12 & $[2,1,2,1,1,2,1,1,1,2,2,2]$ & 2 & 18 & 125761481 \\
13 & $[2,1,2,1,1,2,1,1,1,2,2,2,2]$ & 2 & 20 & 629069549 \\
14 & $[2,1,2,1,1,2,1,1,1,2,2,2,2,1]$ & 1 & 21 & 3146396321 \\
15 & $[2,1,2,1,1,2,1,1,1,2,2,2,2,1,1]$ & 1 & 22 & 15734078757 \\
16 & $[2,1,2,1,1,2,1,1,1,2,2,2,2,1,1,3]$ & 3 & 25 & 78674588089 \\
\bottomrule
\end{tabular}
\end{center}

Every row is an instance of the exact word-orbit identity
\begin{equation*}
  2^{\wordWeight_n}\,\Orbit_7(n)=5^n\cdot7+\wordA_n.
\end{equation*}
For example, the $n=16$ row reads
\begin{equation*}
  2^{25}\cdot34177=5^{16}\cdot7+78674588089.
\end{equation*}
This identity is a theorem in the string-flow layer; it is not a
finite regression fact and does not use \texttt{native\_decide}.

\subsection{Why explicit rewriting}

The finite base is deliberately expanded by explicit rewriting instead
of decision procedures.  A \texttt{native\_decide} certificate would
move the arithmetic into the code-generation trust boundary.  The
explicit form keeps every multiplication and division visible and lets
\texttt{lake build} check the chain without any hidden computation.
The Lean statement
\texttt{fullOrbitIter\_\allowbreak prefix\_\allowbreak expand}
is a conjunction of equalities of the form
\begin{verbatim}
fullOrbitIter 0 = 7 /\ fullOrbitIter 1 = 9 /\
fullOrbitIter 2 = 23 /\ ... /\ fullOrbitIter 16 = 34177
\end{verbatim}

\subsection{Use in the segment exclusions}

The finite base is used in exactly two ways.

First, the $d=3$ survivor produces a candidate first big step of
weight $6$ at depth $j+4$.  By the finite base, the first big step of
the orbit of $7$ occurs at depth $15$ and has weight $3$, not $6$.
The Lean statement
\texttt{candidate\_\allowbreak first\_\allowbreak big\_\allowbreak
step\_\allowbreak weight\_\allowbreak ne\_\allowbreak three} turns this
into a contradiction.

Second, the $d\ge4$ branches produce a candidate first big step of
weight at least $5$, or reduce to the $e=3$, $j=17$ residue
exclusions.  A first big step of weight at least $5$ cannot occur at
depth $15$, where the true weight is $3$, and cannot occur before it,
where all weights are at most $2$; if it occurred later, the true
depth-$15$ step would already be first.  The Lean statement
\texttt{candidate\_\allowbreak first\_\allowbreak big\_\allowbreak
step\_\allowbreak weight\_\allowbreak ge\_\allowbreak five} records this
case.

\subsection{First big step uniqueness}

The uniqueness lemma has the following Lean shape.
\begin{verbatim}
theorem first_big_step_unique (n : Nat)
    (hsmall : forall m : Nat, m < n ->
      twoValuation (5 * fullOrbitIter m + 1) <= 2)
    (hbig : 3 <= twoValuation (5 * fullOrbitIter n + 1)) :
    n = 15 /\ twoValuation (5 * fullOrbitIter n + 1) = 3
\end{verbatim}

\subsection{Related finite facts}

\begin{itemize}[leftmargin=*]
\item \texttt{fullOrbitIter\_\allowbreak prefix\_\allowbreak expand}:
  the seventeen exact values.
\item \texttt{fullOrbit\_\allowbreak prefix\_\allowbreak step\_\allowbreak
  weights}: the exact weights for $n<16$.
\item \texttt{fullOrbit\_\allowbreak first\_\allowbreak t\_\allowbreak
  ge3\_\allowbreak is\_\allowbreak exactly\_\allowbreak 3}: the first
  big step is at depth $15$ with weight $3$.
\item \texttt{first\_\allowbreak big\_\allowbreak step\_\allowbreak
  unique}: uniqueness of the first big step.
\end{itemize}

These four facts are the only output of the finite-prefix chapter that
the unified core consumes.

\subsection{Status}

\begin{center}
\small
\begin{tabular}{@{}ll@{}}
\toprule
Statement & Status \\
\midrule
exact 17-state expansion & Lean: compiled \\
exact step weights & Lean: compiled \\
first big step at depth 15 with weight 3 & Lean: compiled \\
first big step unique & Lean: compiled \\
no \texttt{native\_decide} in the chain & verified by inspection \\
\bottomrule
\end{tabular}
\end{center}
